\section{Theorie}
\label{sec:Theorie}

\subsection{Das Myon}
\label{sec:t1}
Das Myon gehört zu den Elementarteilchen der Leptonen und ist etwa 200 mal 
schwerer als das Elektron. Es hat eine Ruhemasse von $\qty{105.6}{\mega\electronvolt}$
und gehört mit seinem halbzahligen Spin zu der Gruppe der Fermionen. Auffällig ist 
die sehr kurze mittlere Lebensdauer, diese beträgt aufgrund der hohen Masse 
nur $\qty{2.2}{\micro\second}$. \cite{Myon}
Kosmische Myonen entstehen $\sim \qty{30}{\kilo\meter}$ über der Erdoberfläche, 
wo hochenergetische Teilchen, meist Protonen, mit den Atomen der Atmosphäre 
kollidieren. Dabei entstehen Pionen, deren Zerfall wiederum Myonen erzeugen:
\begin{align*}
    \pi^+ \longrightarrow \mu^+ + \nu_\mu, \\
    \pi^- \longrightarrow \mu^- + \nu_\mu
\end{align*}
Aufgrund der Entstehungshöhe und der verhältnismäßig kurzen Lebensdauer dürften
es Myonen klassischerweise nicht viel weiter als $\sim \qty{660}{\meter}$
schaffen. Aufgrund ihrer hohen Energie sind sie sehr schnell, annähernd 
Lichtgeschwindigkeit, wenn sie sich in Richtung des Erdbodens bewegen. Daher
bedarf es relativistischer Betrachtung, welche Aufschluss über
die tatsächliche Reichweite von Myonen gibt:
\begin{equation*}
    s = \gamma v \tau \approx \qty{62.5}{\kilo\meter}, \quad \text{mit } \gamma = \frac{E}{m_\mu c^2}.
\end{equation*}
Der Zerfall von Myonen ergibt
\begin{align*}
    \mu^+ \longrightarrow e^+ + \overline{\nu_e} + \nu_\mu, \\
    \mu^- \longrightarrow e^- + \nu_e + \overline{\nu_\mu},
\end{align*}
also u.a. Elektronen und Positronen. Diese spielen eine zentrale Rolle für 
diesen Versuch. Die Bestimmung der Lebensdauer erfolgt nach der Methode der 
Einzelzerfälle. Ein Myon, das den Szintillatortank (dieser dient als Detektor, 
siehe \autoref{sec:t3}) durchquert, zerfällt unter einer gewissen Wahrscheinlichkeit 
in ein Elektron. Die Signale von Eintritt und Zerfall  werden zeittechnisch 
von der Elektronik erfasst. Aus der Aufzeichnung tausender individueller 
Zerfallszeiten wird an einem angeschlossenen Computer ein Histogramm erstellt, 
dass die Häufigkeit der Zerfälle in Abhängigkeit von Kanälen angibt, welche 
direkten Aufschluss über die Zeitdifferenz der Signale geben (mehr dazu in 
\autoref{t4}). Daraus wird ein Spektrum $N(t)$ erstellt wodurch mithilfe einer 
Ausgleichsrechnung $\tau$ ermittelt wird. 


\subsection{Die Statistik hinter der Lebensdauer des Myons}
\label{sec:t2}
Die Lebensdauer $\tau$ ist definiert über die durchschnittliche Zeit, bis ein 
Teilchen zerfällt. Die Wahrscheinlichkeit, dass das Myon in einer kurzen 
Zeitspanne zerfällt ist proportional zu $\mathrm{d}t$. Der dazugehörige
Proportionalitätsfaktor ist die Zerfallskonstante $\lambda$. 
\begin{equation*}
    \mathrm{d}P = \lambda \cdot \mathrm{d}t
\end{equation*}
Die Abnahme der Teilchenzahl ist über diese Wahrscheinlichkeit definiert:
\begin{equation*}
    \mathrm{d}N = -N \cdot \mathrm{d}P.
\end{equation*}
Aus beiden Differentialen lässt sich die Differentialgleichung 
\begin{equation*}
    \frac{1}{N} \mathrm{d}N = - \lambda \mathrm{d}t
\end{equation*}
aufstellen, dessen Lösung nach Separation der Variablen durch
\begin{equation*}
    N(t) = N_0 e^{-\lambda t}
\end{equation*}
gegeben ist. Dabei ist $N_0=N(0)$ die Anzahl der Myonen zum Anfangszeitpunkt.
Die mittlere Lebensdauer ist der Erwartungswert der Zeit, folglich kann diese
nun als
\begin{equation}
    \tau = \langle t \rangle \\
         = \frac{\int_0^\infty t \cdot N(t) \mathrm{d}t}{\int_0^\infty N(t) \mathrm{d}t} 
         = \frac{\frac{1}{\lambda^2}}{\frac{1}{\lambda}} \\
         = \frac{1}{\lambda}
\end{equation}
bestimmt werden. $\tau$ ist also der Kehrwert der Zerfallskonstanten $\lambda$.
Die Funktion für die Teilchenzahl $N(t)$ lautet infolgedessen
\begin{equation}
    N(t) = N_0 e^{-\frac{t}{\tau}}.
\end{equation}
% Damit lautet die Änderung für die Änderung der Teilchenzahl nun
% \begin{equation*}
%     \mathrm{d}N = - \lambda N_0 e^{-\lambda t} \mathrm{d}t.
% \end{equation*}
% Die Änderung der Zerfälle kann durch $\mathrm{d}N_{\text{Zerfälle}}=-\mathrm{d}N$
% ausgedrückt werden. Die Wahrscheinlichkeitsdichte $P = -\dot{N}{N}$ führt mit


\subsection{Szintillatoren}
\label{sec:t3}
Ein Szintillator ist in der Regel ein Kristall, welcher über die Eigenschaft 
verfügt, beim Durchgang von energiereichen Teilchen durch Stoßprozesse angeregt
zu werden. Diese Anregungsenergie gibt er in Form von Licht wieder ab. Das 
Spektrum umfasst das sichtbare Licht, UV- und Röntgenstrahlen. Es wird zwischen 
zwei Arten von Szintillatoren unterschieden: den organischen und den anorganischen.
Organische Szintillatoren sind etwa Kunststoff, spezielle flüssige Lösungsmittel 
oder Kristalle. Sie haben schnelle Fluoreszenzzeiten (das heißt wenige Nanosekungen)
und eignen sich somit sehr gut für schnelle Messungen. Große Formate für 
präzisere Messungen lassen sich leicht herstellen und sind günstig. Andererseits 
altern diese bei hoher Strahlungsdosis schneller. Darüber hinaus ist die 
Lichtausbeute gegebenenfalls anisotrop. Das bedeutet, dass die Menge an erzeugtem 
Szinitillationslicht abhängig von der Einfallrichtung der zu detektierenden 
Teilchen ist. Anders ist es bei anorganischen Szintillatoren, sie bieten eine 
gute Lichtausbeute und sind dazu noch ziemlich resistent gegenüber Strahlung.
Sie haben jedoch den Nachteil, dass sie wesentlich langsamer als organische 
Szintillatoren sind, ihre Abklingzeit beträgt in der Regel einige hunder Nanosekungen,
was präzise Messungen erschwert. \cite{Szintillator}

\subsection{Fehlerrechnung}
Die gemessenen Werte unterliegen Messunsicherheiten und werden demnach im
Folgenden nicht als fehlerfrei angesehen. Die Fehler entstehen bei der
Bildung der Mittelwerte durch den Fehler des Mittelwerts und bei der
Regressionsrechnung sowie der Fehlerforpflanzung durch Python.
Der Mittelwert ist definiert durch
\begin{equation}
    \overline{x} = \frac{1}{N} \sum\limits_{i=1}^N x_i.
\end{equation}
\noindent Der Fehler des Mittelwerts ist somit gegeben durch 
\begin{equation}
    \begin{aligned}
        \increment \overline{x} &= \sqrt{\overline{x^2\kern-0.1em} - \overline{x}^2} \\
                            &= \frac{\sqrt{\frac{1}{N-1} \sum\limits_{i=1}^N (x_i - \overline{x})^2}}{\sqrt{N}}.
    \end{aligned}
\end{equation}

\noindent Um Fehler einzubeziehen, wird die Gauß'sche Fehlerfortpflanzung verwendet:
\begin{equation}
    \label{eqn:9}
    \increment f = \sqrt{\left(\frac{\partial f}{\partial x}\right)^2 \cdot \left(\increment x\right)^2 + \left(\frac{\partial f}{\partial y}\right)^2 \cdot \left(\increment y\right)^2 + .... + \left(\frac{\partial f}{\partial z}\right)^2 \cdot \left(\increment z\right)^2}
\end{equation}
